% Study-guide LaTeX template — house style for this vault's lecture study
% guides. Copy this, fill in <PLACEHOLDERS>, compile with `tectonic file.tex`.
% See references/tectonic-setup.md for the tectonic install/compile flow.

\documentclass[11pt]{article}
\usepackage[margin=0.9in]{geometry}
\usepackage{amsmath,amssymb}
\usepackage{xcolor}
\usepackage{enumitem}
\usepackage{tcolorbox}
\tcbuselibrary{skins,breakable}
\usepackage{fancyhdr}
\usepackage{titlesec}
\usepackage{parskip}
\usepackage{hyperref}

\definecolor{headingblue}{HTML}{1A3D6D}
\definecolor{subheadblue}{HTML}{2C5A8F}
\definecolor{calloutbg}{HTML}{EEF4FB}
\definecolor{calloutborder}{HTML}{9CC0E8}

\titleformat{\section}{\Large\bfseries\color{headingblue}}{\thesection}{0.6em}{}
\titleformat{\subsection}{\large\bfseries\color{subheadblue}}{\thesubsection}{0.6em}{}
\titleformat{\subsubsection}{\normalsize\bfseries\itshape}{\thesubsubsection}{0.6em}{}

\pagestyle{fancy}
\fancyhf{}
\lfoot{\small\color{gray} <COURSE CODE> Study Guide --- for personal review}
\rfoot{\small\color{gray} Page \thepage}
\renewcommand{\headrulewidth}{0pt}
\renewcommand{\footrulewidth}{0pt}

% Blue-tinted callout box: use for "key insight" / "core fact" boxes.
\newtcolorbox{callout}{colback=calloutbg,colframe=calloutborder,boxrule=0.5pt,arc=2pt,left=6pt,right=6pt,top=4pt,bottom=4pt,breakable}
% White answer box with a thin gray border: use for every answer-key entry.
\newenvironment{answerbox}{\begin{tcolorbox}[colback=white,colframe=gray!50,boxrule=0.4pt,arc=1pt,left=6pt,right=6pt,top=4pt,bottom=4pt,breakable]}{\end{tcolorbox}}

% Shorthand for a bracketed matrix: \vect{1 & 2 \\ 3 & 4}
\newcommand{\vect}[1]{\begin{bmatrix}#1\end{bmatrix}}

% IMPORTANT: never bake "N of M" into the title itself — just "Study Guide N".
% The total-count context belongs in chat/prose, not the artifact (guides get
% added/reordered/split later and a hardcoded total goes stale immediately).
\title{\color{headingblue}\textbf{<COURSE CODE> Study Guide <N>}\\[4pt]\large\normalfont\color{black} <Lecture range>: <Topic summary>}
\author{}
\date{}

\begin{document}
\maketitle
\thispagestyle{fancy}
\vspace{-1.5em}
\noindent\hrulefill

\section*{1. <First topic heading> (Lecture <N>)}
<Walkthrough prose. Ground every claim/example in the actual OCR'd lecture
content — never invent a "similar" example from general knowledge.>

\begin{callout}
\textbf{<Key insight label>:} <the one-sentence takeaway worth a highlighted box.>
\end{callout}

\begin{quote}
Example (Lecture <N>):
\[
\left[\begin{array}{ccc|c} 1 & -2 & 1 & 0 \\ 0 & 1 & -4 & 4 \\ 0 & 0 & 1 & -1\end{array}\right]
\]
<Explain what the matrix shows, worked step by step.>
\end{quote}

% ... repeat \section*{} blocks for each lecture/topic in this guide's range ...

\newpage
\section*{Practice Problems}
<Fresh problems reusing the lecture's exact notation/matrix sizes/difficulty
register — new numbers, same structure, so it's genuine unseen practice.>

\subsubsection*{Problem 1 (<short descriptor>)}
<Problem statement, with a matrix/system in the same \[ \left[\begin{array}...
\end{array}\right] \] style as the walkthrough examples.>

% ... Problem 2, 3, 4, 5 ...

\newpage
\section*{Answer Key}

\subsubsection*{Problem 1}
\begin{answerbox}
<Full worked answer. MUST be independently numerically verified (numpy /
execute_code) before going here — do not hand-wave arithmetic.>
\end{answerbox}

% ... Answer 2, 3, 4, 5, matching problem order ...

\end{document}
